\section{Zielsetzung und Forschungsfragen}
\label{sec:01-03_research-goals}

Ausgehend von der in \hyperref[sec:01-01_problem-and-context]{Abschnitt~1.1 \enquote{Problemstellung und Zweck}} beschriebenen Diskrepanz zwischen dem strategischen Anspruch digitaler Wissensplattformen und ihrer tatsächlichen Nutzung zielt diese Arbeit darauf ab, einen Beitrag zum Verständnis zu leisten, wie \Gls{uxstrategien} und eine systematisch gestaltete \Gls{ia} die Akzeptanz und Wirksamkeit solcher Systeme im organisationalen Kontext verbessern können. Aufbauend auf etablierten Akzeptanzmodellen wie dem \textit{\Gls{tam}} von Davis, das die Bedeutung wahrgenommener Nützlichkeit und Benutzerfreundlichkeit für die Nutzungsintention hervorhebt, sowie dem \textit{\Gls{utaut}} nach Venkatesh et al., das zusätzlich \Gls{performanceexpectancy}, \Gls{effortexpectancy}, sozialen Einfluss und Rahmenbedingungen berücksichtigt, wird \Gls{ux} als zentraler Hebel zur Gestaltung positiver \Gls{nutzungserwartungen} und tatsächlicher Nutzung verstanden \cite{102:Davis1989,103:Venkatesh2003}. In Verbindung mit Ansätzen der digitalen Transformation, wie sie etwa von Westerman et al. diskutiert werden, wird \Gls{ux} nicht nur als Oberflächengestaltung, sondern als strategischer Faktor für die erfolgreiche Verankerung digitaler Arbeitsweisen und Wissensprozesse betrachtet \cite{403:Westerman2014,100:WhartonUX2025}.

Vor diesem theoretischen Hintergrund verfolgt die Arbeit das Ziel, \Gls{uxstrategien} und informationsarchitektonische Gestaltungsprinzipien systematisch zu analysieren und für den Kontext interner Wissens- und Kollaborationsplattformen nutzbar zu machen. Dabei steht insbesondere die Frage im Fokus, wie solche Strategien dazu beitragen können, \gls{pu} und \gls{peou} im Sinne von \Gls{tam} sowie \Gls{performanceexpectancy} und \Gls{effortexpectancy} im Sinne des \Gls{utaut} positiv zu beeinflussen und damit Akzeptanz und Nutzung zu fördern \cite{102:Davis1989,103:Venkatesh2003}. Gleichzeitig soll aufgezeigt werden, welche Bedeutung \Gls{ux}-Methoden wie \Gls{personas}, \Gls{userjourney} oder \Gls{musterkataloge} für die Entwicklung strukturierter, kognitiv entlastender Informationsarchitekturen haben und wie Gestaltungs- und Kommunikationsansätze so kombiniert werden können, dass sie in agilen, komplexen Arbeitsumfeldern tragfähig sind \cite{300:FraunhoferIESE2018,401:Garrett2011,402:Rosenfeld2015}.

Zur Konkretisierung dieser Zielsetzung werden im Rahmen der Arbeit folgende \Glspl{rq} adressiert:

\vspace{1em}
\begin{table}[H]
    \centering
    \begin{tabular}{p{0.1\textwidth} p{0.8\textwidth}}
        \textbf{\Gls{rq}-1} & \textbf{Wie können UX-Strategien systematisch zur Optimierung und Akzeptanz organisationaler Informationsarchitektur beitragen?} \\
        & Diese Frage zielt darauf ab, den Zusammenhang zwischen UX-orientierten Vorgehensmodellen, informationsarchitektonischen Entscheidungen und wahrgenommener Nützlichkeit bzw. Nutzung interner Plattformen im Sinne etablierter Akzeptanzmodelle (TAM, UTAUT) herauszuarbeiten \cite{300:FraunhoferIESE2018,102:Davis1989,103:Venkatesh2003}. \\[1em]

        \textbf{\Gls{rq}-2} & \textbf{Welche Rolle spielen Personas und weitere UX-Methoden bei der Entwicklung nutzerzentrierter Strukturen?} \\
        & Im Zentrum steht die Untersuchung, inwiefern methodische Artefakte wie \Gls{personas}, \Gls{szenarien} oder \Gls{userjourney}  dazu beitragen, \gls{mentalemodelle} und Aufgaben der Nutzergruppen in eine tragfähige \Gls{ia} zu überführen und \gls{cognitiveload} zu reduzieren \cite{401:Garrett2011,402:Rosenfeld2015,500:ISO9241-210}. \\[1em]

        \textbf{\Gls{rq}-3} & \textbf{Welche Gestaltungs- und Kommunikationsansätze fördern die Akzeptanz und Nutzung solcher Plattformen in agilen, komplexen Arbeitsumfeldern?} \\
        & Diese Frage verknüpft die Perspektive nutzerzentrierter Gestaltung mit Erkenntnissen zu Technologieakzeptanz und digitaler Transformation, um zu identifizieren, wie Designentscheidungen, Einführungsstrategien und begleitende Kommunikation zusammenspielen müssen, damit interne Plattformen langfristig in die Arbeitspraktiken eingebettet werden \cite{300:FraunhoferIESE2018,301:Infragistics2015,103:Venkatesh2003,403:Westerman2014}. \\
    \end{tabular}
\end{table}
\vspace{1em}

Die Beantwortung dieser Forschungsfragen soll sowohl zu einer theoretischen Fundierung des Einsatzes von \Gls{uxstrategien} im Kontext organisationaler \Gls{ia} beitragen als auch praxisorientierte Implikationen für die Gestaltung und Einführung interner Wissens- und Kollaborationsplattformen ableiten.